\documentclass[aspectratio=169,11pt]{beamer}
\usepackage{fontspec}
\setsansfont{cmunss.otf}[BoldFont=cmunsx.otf,ItalicFont=cmunsi.otf]
\usepackage{polyglossia}
\setdefaultlanguage{russian}
\setotherlanguage{english}
\newfontfamily\cyrillicfontsf{cmunss.otf}[BoldFont=cmunsx.otf,ItalicFont=cmunsi.otf]
\usepackage{booktabs,tabularx,amsmath}
\usepackage{tikz}
\usetikzlibrary{arrows.meta,positioning,calc}
\usepackage{array}
\definecolor{Ink}{HTML}{182636}
\definecolor{Muted}{HTML}{495D72}
\definecolor{Route}{HTML}{126A63}
\definecolor{Trace}{HTML}{2D5C86}
\definecolor{Rule}{HTML}{CCD8DF}
\definecolor{Paper}{HTML}{FFFFFF}
\setbeamercolor{normal text}{fg=Ink,bg=Paper}
\setbeamercolor{structure}{fg=Route}
\setbeamercolor{frametitle}{fg=Ink}
\setbeamercolor{footline}{fg=Muted}
\setbeamerfont{frametitle}{size=\fontsize{17}{20}\selectfont,series=\bfseries}
\setbeamersize{text margin left=0.75cm,text margin right=0.75cm}
\setbeamertemplate{navigation symbols}{}
\setbeamertemplate{frametitle}{\vspace{0.28cm}\insertframetitle\par\vspace{0.18cm}}
\setbeamertemplate{footline}{\hfill\usebeamercolor[fg]{footline}\fontsize{7}{8}\selectfont\insertframenumber/\inserttotalframenumber\hspace*{0.75cm}\vspace{0.28cm}}
\hypersetup{colorlinks=true,urlcolor=Trace}
\newcommand{\highlight}[1]{{\color{Route}\bfseries #1}}
\newcommand{\muted}[1]{{\color{Muted}#1}}
\newcommand{\slidebody}{\fontsize{10.8}{13.8}\selectfont}
\newcommand{\slidenote}[1]{\par\vspace{0.24cm}{\color{Muted}\fontsize{8.2}{10.5}\selectfont #1\par}}
\newcommand{\legendmark}[1]{\tikz[baseline=-0.6ex]\fill[#1] (0,0) rectangle (0.18,0.12);}
\newcommand{\ce}{\operatorname{CE}}
\newcommand{\relu}{\operatorname{BiReLU}}
\newcommand{\projecttitle}{Компиляция статической маршрутизации\par с нелинейностью до агрегации}
\newcommand{\PresenterName}{}
\newcommand{\PresenterAffiliation}{}
\newcommand{\presenter}{\ifx\PresenterName\empty\else\vspace{0.45cm}\muted{\PresenterName\ifx\PresenterAffiliation\empty\else\par\PresenterAffiliation\fi}\fi}

\newcommand{\maindiagram}{%
\begin{tikzpicture}[x=1cm,y=1cm,
  edge/.style={draw=Muted,-{Latex[length=1.8mm]},line width=0.52pt},
  token/.style={font=\fontsize{9.5}{12}\selectfont,align=center,text=Ink},
  routepart/.style={token,text=Route},
  exactpart/.style={token,text=Trace}]
  \node[token] (x) at (0,0) {Входы\\$x_1,\ldots,x_N$};
  \node[routepart] (stats) at (3.15,0.74) {Общие статистики\\по интервалам};
  \node[exactpart] (select) at (3.15,-0.74) {Входы явных\\маршрутов $S_m$};
  \node[routepart] (tail) at (6.85,0.74) {Сумма BiReLU\\по порогу $-b_m/w_m$};
  \node[exactpart] (correction) at (6.85,-0.74) {Поправка для\\явных маршрутов};
  \node[token] (out) at (10.2,0) {Выход\\$h_m^Q(x)$};
  \draw[edge] (x.east) -- ++(0.33,0) |- (stats.west);
  \draw[edge] (x.east) -- ++(0.33,0) |- (select.west);
  \draw[edge,Route] (stats.east) -- (tail.west);
  \draw[edge,Trace] (select.east) -- (correction.west);
  \draw[edge] (tail.east) -- ++(0.3,0) |- (out.west);
  \draw[edge] (correction.east) -- ++(0.3,0) |- (out.west);
\end{tikzpicture}}

\newcommand{\oneslidecontent}{%
\slidebody
Слой усредняет отклики BiReLU по обученным вероятностям маршрутов.
При выводе компилятор хранит явные маршруты и остаточную массу групп.
\par\vspace{0.4cm}
\centering\maindiagram\par
\vspace{0.46cm}
\begin{columns}[T,onlytextwidth]
\begin{column}{0.31\textwidth}
{\color{Route}\fontsize{25}{28}\selectfont\bfseries $2{,}81\times$}\par
быстрее точного кеша
\end{column}
\begin{column}{0.31\textwidth}
{\color{Route}\fontsize{25}{28}\selectfont\bfseries $10{,}3\times$}\par
меньше явной памяти
\end{column}
\begin{column}{0.31\textwidth}
{\color{Route}\fontsize{25}{28}\selectfont\bfseries $-0{,}35$}\par
п.п. точности на тесте
\end{column}
\end{columns}
\raggedright
\slidenote{База сравнения — исходная маршрутизирующая сеть с точным кешем.
CIFAR-10, точность по трём запускам. Замеры Apple M3 Pro CPU, FP32, batch 1, seed 0.}
}

\hypersetup{pdftitle={M-talk. Компиляция статической маршрутизации с нелинейностью до агрегации}}
\begin{document}

\begin{frame}[plain]
\vspace{0.65cm}
{\fontsize{22}{26}\selectfont\bfseries\projecttitle\par}
\vspace{0.62cm}
{\fontsize{12}{16}\selectfont
Цель исследования — ускорить вывод обученного слоя\par
с индивидуальными порогами нейронов\par
и измерить изменение качества и памяти.\par}
\presenter
\end{frame}

\begin{frame}[t]{Общие статистики для индивидуальных порогов}
\oneslidecontent
\end{frame}

\begin{frame}[t]{Статические маршруты и агрегация хвоста}
\slidebody
\renewcommand{\arraystretch}{1.42}
\begin{tabularx}{\textwidth}{@{}>{\raggedright\arraybackslash}p{4.4cm}X@{}}
\toprule
\textbf{Близкая работа} & \textbf{Используемая идея} \\
\midrule
\href{https://proceedings.mlr.press/v139/tay21a.html}{Synthesizer, 2021} & Статические факторизованные логиты \\
\href{https://arxiv.org/abs/2404.19756}{KAN, 2025} и \href{https://arxiv.org/abs/2602.03879}{TruKAN, 2026} & Одномерные функции перед суммой \\
\href{https://proceedings.mlr.press/v235/ribar24a.html}{SparQ, 2024} и \href{https://arxiv.org/abs/2610.30951}{OpenAttention, 2026} & Агрегация отброшенной части внимания \\
\bottomrule
\end{tabularx}
\par\vspace{0.46cm}
\highlight{Вклад работы}\par
Компиляция сжатого статического слоя. Нейроны с разными порогами
BiReLU читают одни и те же статистики группы.
\end{frame}

\begin{frame}[t]{Слой средних и критерии сравнения}
\slidebody
\begin{columns}[T,onlytextwidth]
\begin{column}{0.58\textwidth}
\highlight{Обучаемый слой}
\[
 h_m^P(x)=\sum_i p_{mi}\,g(w_mx_i+b_m)
\]
\[
 g(u)=\tfrac12u+\tfrac32\max(u,0)
\]
Пара $w_m,b_m$ общая для всех входов нейрона $m$.
Строка $P$ задаёт вероятности входных координат.
\end{column}
\begin{column}{0.36\textwidth}
\highlight{Критерии}\par\vspace{0.15cm}
Точность и кросс-энтропия\par\vspace{0.2cm}
Память массивов и рабочих буферов\par\vspace{0.2cm}
Задержка полного вывода
\end{column}
\end{columns}
\vspace{0.32cm}
\slidenote{Чекпойнты выбираются по валидационной точности. Коэффициенты сжатия подгоняются на обучающем разбиении. Тест служит для итоговой оценки.}
\end{frame}

\begin{frame}[t]{Сжатие вероятностей и точное вычисление}
\slidebody
\begin{columns}[T,onlytextwidth]
\begin{column}{0.46\textwidth}
\highlight{Офлайн}\par\vspace{0.12cm}
Сохраняем top-$K$ маршрутов.\par
Остаточную массу равномерно распределяем по отброшенным входам каждой группы.
\end{column}
\begin{column}{0.46\textwidth}
\highlight{На каждом объекте}\par\vspace{0.12cm}
Считаем число и сумму значений по обе стороны порога. Добавляем поправку для явных маршрутов.
\end{column}
\end{columns}
\vspace{0.23cm}
\[
 \sum_{i\in G} g(w_mx_i+b_m)
 =\tfrac12(w_ms+b_mn)+\tfrac32(w_ms_++b_mn_+).
\]
\slidenote{$n,s$ — число и сумма значений группы. $n_+,s_+$ — те же статистики для $w_mx_i+b_m>0$.}
\vspace{0.15cm}
\highlight{Исполнитель точно вычисляет сжатую модель $Q$ с учётом округления.}\par
Потерю при замене исходных вероятностей $P$ на $Q$ измеряем отдельно.
\end{frame}

\begin{frame}[t]{Полные модели близки по точности на CIFAR-10}
\slidebody
\renewcommand{\arraystretch}{1.36}
\begin{tabularx}{\textwidth}{@{}Xrrr@{}}
\toprule
\textbf{Набор} & \textbf{Скрытых слоёв} & \textbf{Dense} & \textbf{Средние до суммы} \\
\midrule
MNIST & 2 & 98,00 & 95,49 \\
MNIST & 4 & 98,28 & 95,99 \\
CIFAR-10 & 3 & 50,60 & 50,31 \\
CIFAR-10 & 4 & 50,58 & 51,67 \\
\bottomrule
\end{tabularx}
\slidenote{Точность на тесте, \%. Чекпойнты выбраны по валидации, один seed. Бюджеты 500 эпох на MNIST и 2000 на CIFAR-10.}
\end{frame}

\begin{frame}[t]{При ранге 8 оба порядка BiReLU опередили Factor Dense на CIFAR}
\slidebody
\renewcommand{\arraystretch}{1.36}
\begin{tabularx}{\textwidth}{@{}Xrrrr@{}}
\toprule
\textbf{Набор} & \textbf{Эпох} & \textbf{Factor Dense} & \textbf{До суммы} & \textbf{После суммы} \\
\midrule
CIFAR-10 & 800 & 44,60 & \highlight{49,12} & 48,46 \\
Fashion-MNIST & 500 & 85,92 & 85,60 & \highlight{85,94} \\
CIFAR-100 & 400 & 17,73 & \highlight{20,63} & 20,07 \\
\bottomrule
\end{tabularx}
\slidenote{Точность на тесте, \%, среднее по трём запускам. В каждой строке близкое число параметров.
Ранг относится к весам Dense и к логитам маршрутизации соответственно.}
\vspace{0.4cm}
На CIFAR-10 к 1200 эпохам полный Dense дал 50,81\%, компактная модель до суммы — 50,58\%. Число параметров уменьшилось с 570\,570 до 36\,426.
\end{frame}

\begin{frame}[t]{Потеря точности при сжатии — до 0,35 п.п.}
\slidebody
\renewcommand{\arraystretch}{1.33}
\begin{tabularx}{\textwidth}{@{}Xrrr@{}}
\toprule
\textbf{Набор} & \textbf{Точный кеш} & \textbf{Сжатая сеть} & \textbf{Разница, п.п.} \\
\midrule
CIFAR-10 & 50,58\% & 50,23\% & $-0{,}35$ \\
Fashion-MNIST & 85,60\% & 85,46\% & $-0{,}14$ \\
CIFAR-100 & 20,63\% & 20,36\% & $-0{,}27$ \\
\bottomrule
\end{tabularx}
\slidenote{Точность на тесте, три запуска. Бюджет $K=(64,32,64,32)$ выбран на CIFAR-10.
Подгонка коэффициентов и линейной шапки использует 1536 обучающих изображений каждого набора.}
\vspace{0.25cm}
\highlight{Контроль нелинейности на Fashion-MNIST}\par
При одинаковых маршрутах и массах хвоста точная сумма BiReLU дала
85,79\% на валидации, BiReLU от среднего хвоста — 83,43\%.
\end{frame}

\begin{frame}[t]{На CPU сжатая сеть быстрее собственного точного кеша}
\slidebody
\renewcommand{\arraystretch}{1.27}
\begin{tabularx}{\textwidth}{@{}Xrrr@{}}
\toprule
\textbf{Исполнитель FP32} & \textbf{Batch 1, мкс} & \textbf{Batch 128, мс} & \textbf{Память, байт} \\
\midrule
Точный кеш маршрутов & 59,92 & 6,435 & $\approx$2\,288\,760 \\
Сжатый, сортировка & 24,72 & 2,730 & 209\,172 \\
Сжатый, интервалы порогов & \highlight{21,29} & 2,277 & \highlight{221\,460} \\
Dense, Accelerate & 13,78 & 0,198 & 2\,283\,600 \\
\bottomrule
\end{tabularx}
\slidenote{CIFAR-10, seed 0, Apple M3 Pro CPU. Медиана за пакет.
Память включает хранимые массивы и явные рабочие буферы для batch 1.}
\vspace{0.3cm}
Сжатый исполнитель сохранил все 10\,000 предсказаний \emph{сжатой} формулы на тесте для seed 0.
\end{frame}

\begin{frame}[t]{Компиляция полезна для фиксированной маршрутизирующей сети}
\slidebody
\renewcommand{\arraystretch}{1.35}
\begin{tabularx}{\textwidth}{@{}p{3.15cm}>{\raggedright\arraybackslash}X@{}}
\toprule
\textbf{Сравнение} & \textbf{Вывод} \\
\midrule
С точным кешем & В 2,81 раза быстрее и в 10,3 раза меньше явной памяти при batch 1. \\
На трёх наборах & Потеря тестовой точности после сжатия от 0,14 до 0,35 п.п. \\
С Dense & Оптимизированный FP32 Dense быстрее. При близком объёме хранения Factor Dense INT8 точнее на CIFAR-10 и CIFAR-100. \\
\bottomrule
\end{tabularx}
\par\vspace{0.4cm}
\highlight{Дальнейшая проверка}\par
Повторить замеры на другом CPU и проверить слои с несколькими порогами.
\end{frame}
\end{document}
